
I.19 bhava-pratyayo videha-prakiti-layanam
IV.25 visesa-darsina atma-bhava-bhavana-vinivritti

I. Prajna-pada

I.17 vitarka-vicara ananda-asmita-rupa anugamat samprajnata
I.18 virama-pratyayabhyasa-purva samskara-seso ‘nya

Samkhya-Yoga posits Jna as the nondual absolute and Jnavijnana is the
absolute knowing of the absolute. As Schelling notes there must be a point
where knowledge of the absolute and the absolute are one.
To know something, one simply must project it.
Even if one knows something, one knows it again by projecting it again as a
non-knower. Hence the original statement of Avidya in Samkhya-Yoga is as follows:
Absolute Jna in order to be known, must project itself.
The result is the original objectification of Jna: Buddhi tattva

Jna in itself is not a tattva. Its only attribute is "alinga".
It is inconceivable precisely because it is the projection of the original idea,
the pure principle of knowing.

Prakriti-Purusha-Buddhi

Buddhi, the complete specification of Nature in the form of a logical system,
is the initial evolution of Prakriti. The Atman projects Jna and the result
is Jnavijnana.
  Jnavijnana is the Self-production of the original bija.
But this stage is the stage of Buddhi-eva. It is not known through a
prakasa-kriya-sthiti type of ordinary knowing. It is sva-prakasa.

I.19 bhava-pratyayo videha-prakiti-layanam

IV. Dharma-pada

IV.25 visesa-darsina atma-bhava-bhavana-vinivritti
